\documentclass[11pt]{article}

\usepackage[T1]{fontenc}
\usepackage{lmodern}
\usepackage{amsmath,amssymb,amsthm,mathtools}
\usepackage[margin=1.08in]{geometry}
\usepackage{microtype}
\usepackage[colorlinks=true,linkcolor=blue,citecolor=blue,urlcolor=blue]{hyperref}

\newtheorem{theorem}{Theorem}[section]
\newtheorem{proposition}[theorem]{Proposition}
\newtheorem{lemma}[theorem]{Lemma}
\newtheorem{corollary}[theorem]{Corollary}
\theoremstyle{remark}
\newtheorem{remark}[theorem]{Remark}
\newtheorem{problem}[theorem]{Problem}

\newcommand{\C}{\mathbb C}
\newcommand{\D}{\mathbb D}
\newcommand{\Mn}{M_n(\C)}
\newcommand{\alg}{\operatorname{alg}}
\newcommand{\RePart}{\operatorname{Re}}
\newcommand{\diag}{\operatorname{diag}}
\newcommand{\spn}{\operatorname{span}}

\title{Sharp spectral constants for scaled \(q\)-numerical ranges:\\
the quadratic case and an exact completion reduction}
\author{}
\date{\today}

\begin{document}
\maketitle

\begin{abstract}
For \(0<|q|\leq1\), put
\[
 \Omega_q(A)=q^{-1}W_q(A),\qquad
 C_q=\max\left\{1,
 \frac{2|q|}{1+\sqrt{1-|q|^2}}\right\}.
\]
O'Loughlin and Rani conjectured that
\(\|p(A)\|\leq C_q\max_{\Omega_q(A)}|p|\) for every square matrix
\(A\) of order at least two and every polynomial \(p\).
We prove the conjecture, with the sharp constant, for every
\(2\times2\) matrix and, more generally, for every finite-dimensional
matrix whose minimal polynomial has degree at most two.  The result
also holds for finite orthogonal direct sums of such matrices.
For matrices of every size we prove the sharp norm equivalence
\[
 w_q(A)\geq
 \min\left\{|q|,\frac{1+\sqrt{1-|q|^2}}2\right\}\|A\|,
\]
which settles the conjecture for all affine polynomials.
The proof combines the exact \(q\)-ellipse with a two-point
Schur-class extremal calculation and a quadratic-operator canonical
form.  For the unresolved general case we derive the precise
parameterized positive-real completion inequality that an extension
of the sharp numerical-range argument must satisfy.  This separates
the remaining work into two explicit statements: an operator
completion theorem and an algebra-sensitive boundary-measure
construction.  The manuscript deliberately labels those two
statements as open; no claim of a proof of the general conjecture is
made.
\end{abstract}

\section{Statement, normalization, and sharpness}

Throughout, \(n\geq2\).  For a matrix \(A\in M_n(\C)\) and
\(q\in\C\), \(0<|q|\leq1\), define
\[
 W_q(A)=
 \left\{\langle Ax,y\rangle:
 \|x\|=\|y\|=1,\ \langle x,y\rangle=q\right\},
 \qquad
 \Omega_q(A)=\frac1qW_q(A).
\]
Changing the phase of \(y\) shows that
\(\Omega_q(A)=\Omega_{|q|}(A)\).  We therefore write
\[
 r=|q|,\qquad \tau=\sqrt{1-r^2},\qquad
 \kappa=\frac{1+\tau}{r},\qquad
 \rho=\frac2\kappa=\frac{2r}{1+\tau}.
\]
The proposed constant is
\[
 C_r=\max\{1,\rho\}.
\]
Notice that \(\rho=1\) exactly when \(r=4/5\).

\begin{problem}[Scaled \(q\)-numerical-range conjecture]
\label{prob:q-conjecture}
Prove that for every \(A\in M_n(\C)\), \(n\geq2\), and every complex polynomial
\(p\),
\[
 \|p(A)\|\leq C_r\sup_{z\in\Omega_r(A)}|p(z)|.
\]
\end{problem}

Both terms in \(C_r\) are necessary.  The constant polynomial
\(p\equiv1\) forces a constant at least one.  For
\[
 J=\begin{pmatrix}0&1\\0&0\end{pmatrix}
\]
the exact \(q\)-ellipse formula gives
\[
 \Omega_r(J)=\frac{\kappa}{2}\overline{\D}.
\]
Taking \(p(z)=z\) gives the ratio \(2/\kappa=\rho\).

We record the support-function formula that explains the geometry.
Writing \(y=rx+\tau z\), where \(z\perp x\) and \(\|z\|=1\), gives
\begin{equation}
\label{eq:q-support}
 h_{\Omega_r(A)}(e^{i\theta})
 =
 \sup_{\|x\|=1}
 \left[
  \RePart\langle e^{-i\theta}Ax,x\rangle
  +\frac{\tau}{r}
   \left\|P_{x^\perp}(e^{-i\theta}Ax)\right\|
 \right].
\end{equation}
Indeed, after \(x\) is fixed, the phase of \(z\) can be chosen to
make the second summand real and nonnegative.  The \(q\)-numerical
range is convex \cite{LiMehtaRodman1994}; hence comparison of support
functions in \eqref{eq:q-support} gives
\(W(A)\subseteq\Omega_r(A)\).

\begin{proposition}[The normal case]
\label{prop:normal}
Problem~\ref{prob:q-conjecture} holds with constant \(1\) whenever
\(A\) is normal.
\end{proposition}

\begin{proof}
The spectrum of \(A\) is contained in \(W(A)\), hence in
\(\Omega_r(A)\).  The spectral theorem therefore gives
\[
 \|p(A)\|=\max_{\lambda\in\sigma(A)}|p(\lambda)|
 \leq\sup_{z\in\Omega_r(A)}|p(z)|.
\]
\end{proof}

\section{The conjecture for two-by-two matrices}
\label{sec:q-two-by-two}

For \(0<r\leq 1\), set
\[
  \tau=\sqrt{1-r^{2}},\qquad
  \kappa=\kappa(r):=\frac{1+\tau}{r},
\]
and
\[
  C_r:=\max\left\{1,\frac{2}{\kappa}\right\}
  =\max\left\{1,\frac{2r}{1+\sqrt{1-r^{2}}}\right\}.
\]
Thus \(\kappa\geq 1\), with \(\kappa=2\) precisely when \(r=4/5\);
moreover, \(C_r=1\) precisely for \(0<r\leq4/5\).  We write
\[
  \Omega_r(A):=\frac{1}{r}W_r(A).
\]
The standard phase covariance of the \(q\)-numerical range gives
\(\Omega_q(A)=\Omega_{|q|}(A)\), so it is enough to consider real
\(q=r\).

\begin{theorem}
\label{thm:q-two-by-two}
Let \(A\in M_2(\mathbb C)\), let \(0<|q|\leq 1\), and let \(p\) be a
complex polynomial.  Then
\[
  \|p(A)\|
  \leq
  \max\left\{1,
       \frac{2|q|}{1+\sqrt{1-|q|^{2}}}\right\}
  \sup_{z\in\Omega_q(A)}|p(z)|.
\]
For each fixed \(q\), this constant is best possible uniformly over
\(A\in M_2(\mathbb C)\).
\end{theorem}

We first isolate the disk functional-calculus calculation that will be
used in the distinct-eigenvalue case.

\begin{lemma}[A two-point Schur-class calculation]
\label{lem:two-point-schur}
For \(\sigma\geq 0\), let
\[
  B_\sigma=
  \begin{pmatrix}
    1&2\sigma\\
    0&-1
  \end{pmatrix},
  \qquad
  R=R_\sigma:=\sqrt{1+\sigma^{2}}+\sigma.
\]
Then \(\|B_\sigma\|=R\), and for every \(0<a<1\),
\[
  \sup_{\substack{g\in H^\infty(\mathbb D)\\
                  \|g\|_\infty\leq 1}}
       \|g(aB_\sigma)\|
  =\max\{1,aR\}.
\]
\end{lemma}

\begin{proof}
The positive matrix \(B_\sigma^*B_\sigma\) has determinant \(1\) and
trace \(2+4\sigma^2=R^2+R^{-2}\).  Its eigenvalues are therefore
\(R^2\) and \(R^{-2}\), which proves \(\|B_\sigma\|=R\).

The lower bound follows by taking \(g\equiv 1\) and \(g(z)=z\).
Put \(T=aB_\sigma\).  If \(aR\leq 1\), then \(T\) is a contraction,
and von Neumann's inequality gives the upper bound.

Suppose now that
\[
  C:=aR\geq 1.
\]
Since \(B_\sigma^2=I\), if \(u=g(a)\) and \(v=g(-a)\), then
\begin{equation}
\label{eq:g-of-involution}
  g(T)=
  \begin{pmatrix}
    u&\sigma(u-v)\\
    0&v
  \end{pmatrix}.
\end{equation}
The set
\[
  \mathcal K_a
  :=\{(g(a),g(-a)):g\in H^\infty(\mathbb D),\ \|g\|_\infty\leq 1\}
\]
is compact and convex.  By the two-point Pick theorem, its nonconstant
boundary is described by equality in
\[
  \left|\frac{u-v}{1-\overline{u}v}\right|
  \leq \frac{2a}{1+a^2}.
\]
Equality in Schwarz--Pick implies that the corresponding pair is
obtained from a disk automorphism.  If \(|u|=1\) or \(|v|=1\), the
pair instead comes from a unimodular constant.  Consequently every
extreme point of \(\mathcal K_a\) is obtained from either a unimodular
constant or a disk automorphism.  Since the norm in
\eqref{eq:g-of-involution} is a convex function of \((u,v)\), it is
enough to check these functions.

Unimodular constants are immediate.  Up to an irrelevant unimodular
factor, a disk automorphism has the form
\[
  \psi_\zeta(z)=\frac{z-\zeta}{1-\overline{\zeta}z},
  \qquad |\zeta|<1.
\]
Let
\[
  A_\zeta=T-\zeta I,\qquad
  D_\zeta=I-\overline{\zeta}T,
\]
and consider the Hermitian matrix
\[
  \Delta
  :=C^2D_\zeta^*D_\zeta-A_\zeta^*A_\zeta.
\]
Write \(z=|\zeta|^2\) and \(x=\operatorname{Re}\zeta\).  Direct
calculation gives
\begin{align}
  \operatorname{tr}\Delta
  &=
  z(C^4+a^4-2)+C^2-\frac{a^4}{C^2},
  \label{eq:Delta-trace}\\
  \det\Delta
  &=(C^2-1)
  \bigl[
       2(C^2-a^2)(1+a^2)z
       -4a^2(C^2-1)x^2 \notag\\
  &\hspace{42mm}
       -(1-a^4)(C^2+1)z^2
  \bigr].
  \label{eq:Delta-det}
\end{align}
If \(C^4+a^4-2\geq 0\), then \eqref{eq:Delta-trace} is nonnegative
by setting \(z=0\).  If \(C^4+a^4-2<0\), then \(0\leq z<1\) gives
\[
  \operatorname{tr}\Delta
  \geq
  C^4+a^4-2+C^2-\frac{a^4}{C^2}
  =
  \frac{(C^2-1)(C^4+2C^2+a^4)}{C^2}
  \geq 0.
\]
Since \(x^2\leq z\) and \(z^2\leq z\), the expression in square
brackets in \eqref{eq:Delta-det} is bounded below by
\begin{align*}
  z\bigl[
       2(C^2-a^2)(1+a^2)
       -4a^2(C^2-1)
       -(1-a^4)(C^2+1)
    \bigr]
  =z(C^2-1)(1-a^2)^2.
\end{align*}
Thus \(\operatorname{tr}\Delta\geq 0\) and
\(\det\Delta\geq 0\), so \(\Delta\succeq 0\).  It follows that
\[
  \|\psi_\zeta(T)\|
  =\|A_\zeta D_\zeta^{-1}\|
  \leq C.
\]
This proves the required upper bound.
\end{proof}

\begin{remark}[Algebra audit]
\label{rem:algebra-audit}
The identities \eqref{eq:Delta-trace} and \eqref{eq:Delta-det}, as
well as both displayed factorizations used to prove their
nonnegativity, were independently checked in exact symbolic algebra.
They were also spot-checked numerically.  No algebraic identity in
Lemma~\ref{lem:two-point-schur} remains unchecked.
\end{remark}

We shall also need the coalescing-eigenvalue version.

\begin{lemma}[The square-zero limit]
\label{lem:square-zero-schur}
Let
\[
  J=
  \begin{pmatrix}
    0&1\\
    0&0
  \end{pmatrix}.
\]
For every \(c\geq 0\),
\[
  \sup_{\substack{g\in H^\infty(\mathbb D)\\
                  \|g\|_\infty\leq 1}}
       \|g(cJ)\|
  =\max\{1,c\}.
\]
\end{lemma}

\begin{proof}
Again, \(g\equiv1\) and \(g(z)=z\) give the lower bound.  If
\(c\leq1\), the upper bound is von Neumann's inequality.  Suppose
\(c\geq1\).  Since \(J^2=0\),
\[
  g(cJ)=
  \begin{pmatrix}
    g(0)&c\,g'(0)\\
    0&g(0)
  \end{pmatrix}.
\]
Put \(b=|g(0)|\).  Schwarz--Pick gives
\(|g'(0)|\leq1-b^2\).  A \(2\times2\) upper triangular matrix
\(\left(\begin{smallmatrix}b&d\\0&b\end{smallmatrix}\right)\)
with \(b\leq1\leq c\) has norm at most \(c\) if and only if
\[
  (c^2-b^2)^2\geq c^2|d|^2.
\]
Here \(|d|\leq c(1-b^2)\), and
\[
  c^2-b^2-c^2(1-b^2)=b^2(c^2-1)\geq0.
\]
Hence \(\|g(cJ)\|\leq c\).
\end{proof}

\begin{proof}[Proof of Theorem~\ref{thm:q-two-by-two}]
The affine covariance
\[
  \Omega_r(\alpha A+\beta I)
  =\alpha\Omega_r(A)+\beta
  \qquad(\alpha,\beta\in\mathbb C)
\]
shows that the assertion is invariant under affine changes of the
matrix.

Assume first that \(A\) has two distinct eigenvalues.  Schur
triangularization, followed by an affine change and a diagonal unitary
similarity, reduces \(A\) to
\[
  B_\sigma=
  \begin{pmatrix}
    1&2\sigma\\
    0&-1
  \end{pmatrix},
  \qquad \sigma\geq0.
\]
Put
\[
  v=\sqrt{1+\sigma^2},\qquad R=v+\sigma=\|B_\sigma\|.
\]
The Hermitian and skew-Hermitian parts of \(B_\sigma\) are
\[
  H=
  \begin{pmatrix}
    1&\sigma\\
    \sigma&-1
  \end{pmatrix},
  \qquad
  K=
  \begin{pmatrix}
    0&-i\sigma\\
    i\sigma&0
  \end{pmatrix}.
\]
Thus
\[
  \operatorname{tr}(H^2)=2(1+\sigma^2),\qquad
  \operatorname{tr}(K^2)=2\sigma^2,\qquad
  \operatorname{tr}(HK)=0.
\]
The two-by-two \(q\)-elliptical range formula
\cite[Lemma~6]{ChienNakazatoPsarrakos2005} therefore shows that
\(\Omega_r(B_\sigma)\) is the ellipse centered at the origin with
semiaxes
\begin{equation}
\label{eq:q-ellipse-axes}
  A_r=\frac{v+\tau\sigma}{r},
  \qquad
  B_r=\frac{\sigma+\tau v}{r}.
\end{equation}
These axes satisfy
\[
  A_r^2-B_r^2=1,
  \qquad
  A_r+B_r=\kappa(v+\sigma)=\kappa R.
\]
Consequently \(\Omega_r(B_\sigma)=E_\eta\), where
\[
  \eta:=\kappa R
\]
and \(E_\eta\) is the confocal ellipse
\[
  E_\eta=
  \left\{x+iy:
  \frac{x^2}{((\eta+\eta^{-1})/2)^2}
  +
  \frac{y^2}{((\eta-\eta^{-1})/2)^2}
  \leq1
  \right\}.
\]

Suppose first that \(\eta>1\).  Let
\[
  \phi_\eta:E_\eta^\circ\longrightarrow\mathbb D
\]
be normalized by \(\phi_\eta(0)=0\) and
\(\phi_\eta'(0)>0\).  Symmetry and uniqueness imply that
\(\phi_\eta\) is odd.  Write
\[
  a_\eta:=\phi_\eta(1)\in(0,1).
\]
We claim that
\begin{equation}
\label{eq:focus-map-bound}
  a_\eta\leq\frac{2}{\eta}.
\end{equation}
For completeness, let \(\xi=\log\eta\).  The classical Schwarz map in
the normalization of Kanas--Sugawa
\cite[Equation~(2.1) and the inverse formula following it]
{KanasSugawa2006} chooses \(m\in(0,1)\) so that
\[
  \frac{\pi}{2}\frac{K(m')}{K(m)}=2\xi
\]
where we use the modulus convention
\[
 K(m)=\int_0^{\pi/2}\frac{dt}{\sqrt{1-m^2\sin^2t}},
 \qquad m'=\sqrt{1-m^2},
\]
and gives
\[
  \phi_\eta(w)
  =
  \sqrt m\,
  \operatorname{sn}\left(
     \frac{2K(m)}{\pi}\arcsin w,m
  \right).
\]
In particular, \(a_\eta=\sqrt m\).  If \(Q\) denotes the elliptic
nome, then
\[
  Q=\exp\left(-\pi\frac{K(m')}{K(m)}\right)
   =e^{-4\xi}=\eta^{-4}
\]
and the standard nome identities
\cite[Equations~22.2.1--22.2.2]{DLMF}
give
\[
  m=\frac{\vartheta_2(0,Q)^2}{\vartheta_3(0,Q)^2},
  \qquad
  a_\eta
  =\frac{\vartheta_2(0,Q)}{\vartheta_3(0,Q)}.
\]
Using the theta series
\cite[Equations~20.2.2--20.2.3]{DLMF},
\[
  \vartheta_2(0,Q)
  =2Q^{1/4}\sum_{n=0}^{\infty}Q^{n(n+1)},
  \qquad
  \vartheta_3(0,Q)
  =1+2\sum_{n=1}^{\infty}Q^{n^2}.
\]
Since \(0<Q<1\) and \(n(n+1)\geq n^2\),
\[
  \sum_{n=0}^{\infty}Q^{n(n+1)}
  \leq
  \sum_{n=0}^{\infty}Q^{n^2}
  <\vartheta_3(0,Q).
\]
As \(Q^{1/4}=\eta^{-1}\), this proves
\eqref{eq:focus-map-bound}.

Because \(B_\sigma^2=I\) and \(\phi_\eta\) is odd,
\[
  \phi_\eta(B_\sigma)=a_\eta B_\sigma.
\]
For a polynomial \(p\), define
\[
  g=p\circ\phi_\eta^{-1}.
\]
The conformal map extends continuously to the boundary of the ellipse,
so \(g\in H^\infty(\mathbb D)\) and
\(\|g\|_{\infty,\mathbb D}=\|p\|_{\infty,E_\eta}\).
The composition rule for the holomorphic functional calculus and
Lemma~\ref{lem:two-point-schur} now give
\begin{align*}
  \|p(B_\sigma)\|
  &=\|g(a_\eta B_\sigma)\|\\
  &\leq
  \max\{1,a_\eta R\}\,
  \|p\|_{\infty,E_\eta}\\
  &\leq
  \max\left\{1,\frac{2R}{\eta}\right\}
  \|p\|_{\infty,E_\eta}\\
  &=
  \max\left\{1,\frac{2}{\kappa}\right\}
  \|p\|_{\infty,\Omega_r(B_\sigma)}.
\end{align*}

The only distinct-eigenvalue case with \(\eta=1\) is
\(r=1\) and \(\sigma=0\).  Then \(B_\sigma\) is normal and
\(\Omega_1(B_\sigma)=[-1,1]\), so the spectral theorem gives the
stronger constant \(1\).

Suppose next that \(A\) has a repeated eigenvalue.  If \(A\) is scalar,
the assertion is immediate.  Otherwise, affine covariance reduces the
problem to
\[
  J=
  \begin{pmatrix}
    0&1\\
    0&0
  \end{pmatrix}.
\]
The same two-by-two \(q\)-elliptical range formula gives
\[
  \Omega_r(J)=
  \left\{z\in\mathbb C:|z|\leq\frac{\kappa}{2}\right\}.
\]
Put \(c=2/\kappa\) and, for a polynomial \(p\), set
\(g(z)=p((\kappa/2)z)\).  Then
\[
  p(J)=g(cJ),\qquad
  \|g\|_{\infty,\mathbb D}
  =\|p\|_{\infty,\Omega_r(J)}.
\]
Lemma~\ref{lem:square-zero-schur} yields
\[
  \|p(J)\|
  \leq
  \max\left\{1,\frac{2}{\kappa}\right\}
  \|p\|_{\infty,\Omega_r(J)}.
\]
This completes the proof of the inequality.

Finally, for \(A=J\), the constant polynomial \(p\equiv1\) forces every
spectral constant to be at least \(1\), while \(p(z)=z\) gives the
ratio
\[
  \frac{\|J\|}
       {\sup_{z\in\Omega_r(J)}|z|}
  =\frac{2}{\kappa}.
\]
Hence the best uniform constant over \(M_2(\mathbb C)\) is exactly
\(C_r\).
\end{proof}

 
\section{Quadratic matrices in arbitrary finite dimension}
\label{sec:q-quadratic}

The two-by-two result extends without loss to every finite-dimensional
matrix whose minimal polynomial has degree at most two.  In fact, both
the scaled \(q\)-numerical range and every polynomial norm are already
determined by one extremal \(2\times2\) block.

\begin{proposition}[Exact reduction to an extremal block]
\label{prop:quadratic-reduction}
Let \(A\in M_n(\mathbb C)\) be nonscalar and suppose that
\[
  (A-aI)(A-bI)=0
\]
for some \(a,b\in\mathbb C\), where \(a=b\) is allowed.  Then there is
a number \(d\geq0\) such that, for
\[
  A_0=
  \begin{pmatrix}
    a&d\\
    0&b
  \end{pmatrix},
\]
the following identities hold for every \(0<|q|\leq1\) and every
complex polynomial \(f\):
\[
  \Omega_q(A)=\Omega_q(A_0),
  \qquad
  \|f(A)\|=\|f(A_0)\|.
\]
\end{proposition}

\begin{proof}
The canonical form for a quadratic operator
\cite[Theorem~1.1]{LiPoonSze2011} gives an orthogonal decomposition
\[
  \mathbb C^n=\mathcal H_1\oplus\mathcal H_1\oplus\mathcal H_2
\]
with respect to which
\begin{equation}
\label{eq:quadratic-block-model}
  A\ \simeq\
  \begin{pmatrix}
    aI_{\mathcal H_1}&P\\
    0&bI_{\mathcal H_1}
  \end{pmatrix}
  \oplus \gamma I_{\mathcal H_2}.
\end{equation}
Here \(P\) is positive semidefinite, \(\gamma\in\{a,b\}\), and
\(\mathcal H_1\neq\{0\}\).  When \(a=b\), \(P\) has zero kernel.  Set
\[
  d:=\|P\|.
\]
Since the space is finite dimensional, \(P\) has a unit eigenvector
for the eigenvalue \(d\).  The quadratic-operator elliptical range
theorem
\cite[Theorem~3.1 and Corollary~3.3]{LiPoonSze2011} therefore gives
\[
  W_r(A)=W_r(A_0)
  \qquad(0\leq r\leq1).
\]
After division by \(r\), and then use of the standard phase covariance
\(\Omega_q=\Omega_{|q|}\), this proves
\(\Omega_q(A)=\Omega_q(A_0)\) for \(0<|q|\leq1\).

It remains to prove the norm identity.  Let
\[
  d=d_1\geq d_2\geq\cdots\geq d_m\geq0
\]
be the eigenvalues of \(P\), counted with multiplicity, where
\(m=\dim\mathcal H_1\).  Diagonalizing \(P\) in both copies of
\(\mathcal H_1\) in \eqref{eq:quadratic-block-model} shows that
\begin{equation}
\label{eq:quadratic-direct-sum}
  A\ \simeq\
  A_{d_1}\oplus\cdots\oplus A_{d_m}
  \oplus\gamma I_{\mathcal H_2},
  \qquad
  A_t:=
  \begin{pmatrix}
    a&t\\
    0&b
  \end{pmatrix}.
\end{equation}

For a polynomial \(f\), put
\[
  \delta_f:=
  \begin{cases}
    \displaystyle\frac{f(a)-f(b)}{a-b},&a\neq b,\\[3mm]
    f'(a),&a=b.
  \end{cases}
\]
The polynomial functional calculus gives, in both cases,
\begin{equation}
\label{eq:polynomial-on-quadratic-block}
  f(A_t)=
  \begin{pmatrix}
    f(a)&t\,\delta_f\\
    0&f(b)
  \end{pmatrix}.
\end{equation}
For fixed \(u,v,w\in\mathbb C\), define
\[
  M_t=
  \begin{pmatrix}
    u&tw\\
    0&v
  \end{pmatrix},
  \qquad t\geq0.
\]
If
\[
  S_t:=|u|^2+|v|^2+t^2|w|^2,
\]
then the trace and determinant of \(M_t^*M_t\) are \(S_t\) and
\(|uv|^2\), respectively.  Hence
\begin{equation}
\label{eq:triangular-norm-formula}
  \|M_t\|^2
  =
  \frac{S_t+\sqrt{S_t^2-4|uv|^2}}{2}.
\end{equation}
The right-hand side is nondecreasing in \(S_t\), and \(S_t\) is
nondecreasing in \(t\).  Thus
\[
  0\leq t_1\leq t_2
  \quad\Longrightarrow\quad
  \|M_{t_1}\|\leq\|M_{t_2}\|.
\]
Applying this observation to
\eqref{eq:polynomial-on-quadratic-block} yields
\[
  \|f(A_{d_j})\|\leq\|f(A_d)\|
  =\|f(A_0)\|
  \qquad(1\leq j\leq m).
\]
Moreover, since \(\gamma\in\{a,b\}\),
\[
  |f(\gamma)|
  \leq r\bigl(f(A_0)\bigr)
  \leq\|f(A_0)\|,
\]
where \(r(\,\cdot\,)\) denotes the spectral radius.  Taking norms in
\eqref{eq:quadratic-direct-sum} now gives
\[
  \|f(A)\|
  =
  \max\left\{
       \max_{1\leq j\leq m}\|f(A_{d_j})\|,
       |f(\gamma)|
  \right\}
  =\|f(A_0)\|,
\]
as required.
\end{proof}

\begin{theorem}[The conjecture for quadratic matrices]
\label{thm:q-quadratic}
Let \(A\in M_n(\mathbb C)\), \(n\geq2\), have minimal polynomial of degree at most
two.  Then, for every \(0<|q|\leq1\) and every complex polynomial
\(p\),
\[
  \|p(A)\|
  \leq
  \max\left\{1,
       \frac{2|q|}{1+\sqrt{1-|q|^2}}\right\}
  \sup_{z\in\Omega_q(A)}|p(z)|.
\]
For every \(n\geq2\) and every fixed \(q\), this constant is best
possible uniformly over the quadratic matrices in \(M_n(\mathbb C)\).
\end{theorem}

\begin{proof}
If the minimal polynomial has degree one, then \(A\) is scalar and
\[
  \|p(A)\|=|p(A)|
  =\sup_{z\in\Omega_q(A)}|p(z)|.
\]
Suppose that the minimal polynomial has degree two.  It either has two
distinct roots, in which case \(A\) is annihilated by
\((z-a)(z-b)\) with \(a\neq b\), or it has a repeated root, in which
case \((A-aI)^2=0\).  Both cases are covered by
Proposition~\ref{prop:quadratic-reduction}.  Thus there is a
\(2\times2\) matrix \(A_0\) such that
\[
  \Omega_q(A)=\Omega_q(A_0),
  \qquad
  \|p(A)\|=\|p(A_0)\|.
\]
The asserted inequality follows directly from
Theorem~\ref{thm:q-two-by-two}.

This argument also covers all degeneracies of the extremal block.  If
the two roots are distinct and \(d=0\), then \(A_0\) is normal.  If
the roots coincide, then \(A_0-aI\) is square-zero; the nonscalar
assumption forces \(d>0\).  Scalar matrices were handled separately.

To see sharpness in every dimension \(n\geq2\), take
\[
  A=J\oplus0_{n-2},
  \qquad
  J=
  \begin{pmatrix}
    0&1\\
    0&0
  \end{pmatrix}.
\]
This matrix is quadratic, and
Proposition~\ref{prop:quadratic-reduction} reduces it to \(J\) itself.
As in the proof of Theorem~\ref{thm:q-two-by-two}, the constant
polynomial gives the lower bound \(1\), while \(p(z)=z\) gives the
lower bound
\[
  \frac{2|q|}{1+\sqrt{1-|q|^2}}.
\]
\end{proof}

 
\section{Reducing direct sums}

The preceding result is stable under finite orthogonal direct
sums.  This gives a class considerably larger than the matrices with
quadratic minimal polynomial.

\begin{theorem}[Direct sums of quadratic matrices]
\label{thm:direct-sums}
Suppose
\[
 A=A_1\oplus\cdots\oplus A_m
\]
and every \(A_j\) has minimal polynomial of degree at most two.
Then Problem~\ref{prob:q-conjecture} holds for \(A\), with the sharp
uniform constant \(C_r\).
\end{theorem}

\begin{proof}
If the \(j\)-th summand has dimension at least two, vectors supported
in that summand show that
\[
 W_r(A_j)\subseteq W_r(A),
 \qquad
 \Omega_r(A_j)\subseteq\Omega_r(A).
\]
Theorem~\ref{thm:q-quadratic} then gives
\[
 \|p(A_j)\|
 \leq C_r\sup_{\Omega_r(A_j)}|p|
 \leq C_r\sup_{\Omega_r(A)}|p|.
\]
If \(A_j=[\lambda]\) is one dimensional, then
\(\lambda\in W(A)\subseteq\Omega_r(A)\), and therefore
\[
 \|p(A_j)\|=|p(\lambda)|
 \leq\sup_{\Omega_r(A)}|p|.
\]
Taking the maximum over \(j\), and using
\(\|p(A)\|=\max_j\|p(A_j)\|\), proves the assertion.  Sharpness
already occurs for a summand equal to \(J\).
\end{proof}

\section{The linear-polynomial checkpoint}
\label{sec:linear-checkpoint}

Testing Problem~\ref{prob:q-conjecture} with \(p(z)=z\) suggests the
inequality
\begin{equation}
\label{eq:linear-checkpoint}
 w_r(A)\geq
 \min\left\{r,\frac{1+\sqrt{1-r^2}}2\right\}\|A\|,
 \qquad
 w_r(A):=\sup_{z\in W_r(A)}|z|.
\end{equation}
The constants are forced respectively by scalar matrices and by
rank-one nilpotents with orthogonal initial and final vectors.

We first prove the scalar inequality underlying the result.  A
phase and a unitary similarity put an arbitrary \(2\times2\) matrix
in the Pauli form
\[
 A=(u+iv)I+h\sigma_z+ik\sigma_x,
 \qquad h\geq k\geq0.
\]
To see this, subtract half the trace and write the traceless part as
\(H+iK\), with \(H,K\) traceless Hermitian.  A scalar phase rotates
the real two-plane spanned by \(H,K\); choose it so that the rotated
pair is Hilbert--Schmidt orthogonal and order the pair by norm.
The Pauli-vector identification of the traceless Hermitian matrices
with \(\mathbb R^3\), together with the \(SO(3)\) action induced by
unitary similarity, then sends the pair to
\(h\sigma_z,k\sigma_x\).
Writing \(X=|u|\), \(Y=|v|\), one has
\begin{align}
\label{eq:pauli-norm}
 \|A\|^2
 &=
 X^2+Y^2+h^2+k^2
 +2\sqrt{h^2X^2+k^2Y^2+h^2k^2},\\
\label{eq:pauli-lower}
 w_r(A)^2
 &\geq
 \max\left\{
 r^2Y^2+(rX+h+\tau k)^2,\,
 r^2X^2+(rY+k+\tau h)^2
 \right\}.
\end{align}
The first identity follows by computing the larger eigenvalue of
\(A^*A\).  The second follows from the exact two-by-two
\(q\)-ellipse formula \cite[Lemma~6]{ChienNakazatoPsarrakos2005}.
The two quantities in \eqref{eq:pauli-lower} are the squared moduli
of the horizontal and vertical extreme points of the exact
\(q\)-ellipse.  Hence the sharp linear checkpoint is implied by the
elementary scalar inequality
\begin{align}
\label{eq:scalar-checkpoint}
 &\max\left\{
 r^2Y^2+(rX+h+\tau k)^2,\,
 r^2X^2+(rY+k+\tau h)^2
 \right\}\\
 &\qquad\geq
 \min\left\{r^2,\frac{(1+\tau)^2}{4}\right\}
 \left(
 X^2+Y^2+h^2+k^2+
 2\sqrt{h^2X^2+k^2Y^2+h^2k^2}
 \right).\notag
\end{align}
This reduction is useful independently of the nonlinear functional
calculus: any proposed proof of the full conjecture must, at a
minimum, account for the phase transition already visible in
\eqref{eq:scalar-checkpoint}.

\begin{lemma}[The sharp scalar inequality]
\label{lem:sharp-scalar-q}
Let \(0<r\leq 1\), put
\[
 s=(1-r^2)^{1/2},
 \qquad
 c=\min\left\{r,\frac{1+s}{2}\right\},
\]
and suppose that \(X,Y\geq0\) and \(h\geq k\geq0\).  Define
\[
 \begin{split}
 N^2&=X^2+Y^2+h^2+k^2
      +2\sqrt{h^2X^2+k^2Y^2+h^2k^2},\\
 M_1^2&=r^2Y^2+(rX+h+sk)^2,\\
 M_2^2&=r^2X^2+(rY+k+sh)^2.
 \end{split}
\]
Then
\[
 \max\{M_1,M_2\}\geq cN.
\]
The constant is sharp: if \(c=r\), equality is attained by
\(h=k=0\) and \((X,Y)\ne(0,0)\); if
\(c=(1+s)/2\), equality is attained by \(X=Y=0\) and
\(h=k>0\).
\end{lemma}

\begin{proof}
Set
\[
 \alpha=h+sk,\qquad \beta=k+sh.
\]
We first show that there is a number \(\lambda\in[0,1]\) such that
\begin{equation}
\label{eq:lambda-constraints}
 \lambda\alpha\geq ch,
 \qquad
 (1-\lambda)\beta\geq ck.
\end{equation}
If \(h=k=0\), this is immediate.  Suppose \(h>0\), and write
\(t=k/h\in[0,1]\).  When a quotient below has both numerator and
denominator equal to zero, it is to be read as zero.  We have
\[
 \frac{h}{\alpha}+\frac{k}{\beta}
   =\frac{1}{1+st}+\frac{t}{t+s}.
\]
For \(s>0\), the derivative of the right-hand side is
\[
 s\left(\frac{1}{(t+s)^2}
          -\frac{1}{(1+st)^2}\right)\geq0,
\]
because
\[
 t+s\leq1+st
 \quad\Longleftrightarrow\quad
 (1-s)(1-t)\geq0.
\]
Consequently,
\[
 \frac{h}{\alpha}+\frac{k}{\beta}
 \leq\frac{2}{1+s}.
\]
The same estimate holds when \(s=0\): the left-hand side is \(2\)
for \(t>0\), and is \(1\) under the above convention when \(t=0\).
Since \(c\leq(1+s)/2\), it follows that
\[
 \frac{ch}{\alpha}+\frac{ck}{\beta}\leq1.
\]
Thus the interval
\[
 \left[\frac{ch}{\alpha},\,1-\frac{ck}{\beta}\right]
\]
is nonempty, and any \(\lambda\) in this interval satisfies
\eqref{eq:lambda-constraints}.

Using
\[
 \begin{split}
 M_1^2&=r^2(X^2+Y^2)+\alpha^2+2rX\alpha,\\
 M_2^2&=r^2(X^2+Y^2)+\beta^2+2rY\beta,
 \end{split}
\]
we obtain
\begin{align*}
 \max\{M_1^2,M_2^2\}
 &\geq \lambda M_1^2+(1-\lambda)M_2^2\\
 &=r^2(X^2+Y^2)
   +\lambda\alpha^2+(1-\lambda)\beta^2\\
 &\hspace{2cm}
   +2r\bigl(\lambda X\alpha+(1-\lambda)Y\beta\bigr).
\end{align*}
Now \(c\leq r\), and \eqref{eq:lambda-constraints} gives
\[
 r\lambda\alpha\geq rch\geq c^2h,
 \qquad
 r(1-\lambda)\beta\geq rck\geq c^2k.
\]
Moreover, by convexity of the square and
\eqref{eq:lambda-constraints},
\[
 \lambda\alpha^2+(1-\lambda)\beta^2
 \geq\bigl(\lambda\alpha+(1-\lambda)\beta\bigr)^2
 \geq c^2(h+k)^2.
\]
It follows that
\begin{equation}
\label{eq:strong-scalar-majorant}
 \max\{M_1^2,M_2^2\}
 \geq c^2\bigl(
 X^2+Y^2+(h+k)^2+2hX+2kY
 \bigr).
\end{equation}
Finally, all three quantities \(hX,kY,hk\) are nonnegative, so
\[
 \sqrt{h^2X^2+k^2Y^2+h^2k^2}
 \leq hX+kY+hk.
\]
Therefore
\[
 N^2
 \leq X^2+Y^2+h^2+k^2+2hX+2kY+2hk
 =X^2+Y^2+(h+k)^2+2hX+2kY.
\]
Together with \eqref{eq:strong-scalar-majorant}, this proves
\(\max\{M_1,M_2\}\geq cN\).

For sharpness, if \(h=k=0\), then
\[
 N=(X^2+Y^2)^{1/2},
 \qquad
 M_1=M_2=r(X^2+Y^2)^{1/2}.
\]
Thus equality holds whenever \(c=r\).  On the other hand, if
\(X=Y=0\) and \(h=k>0\), then
\[
 N=2h,
 \qquad
 M_1=M_2=(1+s)h,
\]
so equality holds whenever \(c=(1+s)/2\).
\end{proof}
 
\begin{theorem}[Sharp norm equivalence for the \(q\)-numerical radius]
\label{thm:q-radius-norm}
For every \(A\in M_n(\C)\), \(n\geq2\), and every \(0<r\leq1\),
\[
 w_r(A)\geq
 \min\left\{r,\frac{1+\sqrt{1-r^2}}2\right\}\|A\|.
\]
The constant is best possible in every dimension at least two.
\end{theorem}

\begin{proof}
We first prove the assertion for \(B\in M_2(\C)\).  Multiplication by
a scalar phase and unitary similarity put \(B\) in the Pauli form
used above.  The exact norm formula is \eqref{eq:pauli-norm}.
The \(q\)-ellipse of \(B\) is centered at \(r(u+iv)\), has a
horizontal semiaxis \(h+\tau k\), and a vertical semiaxis
\(k+\tau h\).  Its horizontal and vertical extreme points give
\eqref{eq:pauli-lower}.  Lemma~\ref{lem:sharp-scalar-q}, with
\[
 c=\min\left\{r,\frac{1+\tau}{2}\right\},
\]
therefore yields \(w_r(B)\geq c\|B\|\).

For a general \(A\in M_n(\C)\), \(n\geq2\), choose a unit right singular vector
\(x\) such that \(\|Ax\|=\|A\|\), and put
\(\mathcal E=\spn\{x,Ax\}\).  If \(\dim\mathcal E=1\), enlarge
\(\mathcal E\) to any two-dimensional subspace.  The compression
\[
 B=P_{\mathcal E}A|_{\mathcal E}
\]
satisfies
\[
 \|B\|\geq\|Bx\|=\|Ax\|=\|A\|,
\]
while the reverse inequality is automatic; hence \(\|B\|=\|A\|\).
Every admissible pair of vectors in \(\mathcal E\) is also admissible
in the full space, so \(W_r(B)\subseteq W_r(A)\).  Consequently
\[
 w_r(A)\geq w_r(B)\geq c\|B\|=c\|A\|.
\]

Scalar matrices show sharpness of the branch \(c=r\).  For the other
branch the \(q\)-ellipse calculation gives
\(w_r(J)=(1+\tau)/2\).  In higher dimensions take
\(A=J\oplus0\).  Proposition~\ref{prop:quadratic-reduction} gives
\(W_r(A)=W_r(J)\), so these are sharp examples in every dimension at
least two.
\end{proof}

Since
\[
 \sup_{z\in\Omega_r(A)}|z|=\frac{w_r(A)}r,
\]
Theorem~\ref{thm:q-radius-norm} proves
Problem~\ref{prob:q-conjecture} for \(p(z)=z\), with the sharp
constant \(C_r=r/c\).  More generally, for
\(p(z)=\alpha z+\beta\), affine covariance gives
\[
 \Omega_r(p(A))=p(\Omega_r(A)).
\]
Applying Theorem~\ref{thm:q-radius-norm} to \(p(A)\) therefore proves
the conjecture for every affine polynomial, again with \(C_r\).
This is only a checkpoint for the nonlinear functional calculus:
\(q\)-numerical ranges do not in general obey a polynomial
spectral-mapping inclusion.

\section{A parameterized positive-real completion}
\label{sec:parameterized-completion}

The proof that the numerical range is a \(2\)-spectral set uses a
positive-real completion of the resolvent.  The correct interpolation
of that completion for the present problem is
\[
 F_\rho(w,T)
 =
 \left(1-\frac2\rho\right)I+
 \frac2\rho(I-wT)^{-1},
 \qquad 0<\rho\leq2.
\]
Without an adjoint-algebra correction, the condition
\(\RePart F_\rho(w,T)\succeq0\) is the standard Herglotz
characterization of a \(\rho\)-contraction.

\begin{proposition}[A necessary Gramian inequality]
\label{prop:rho-gramian}
Let \(T\in M_n(\C)\) have distinct eigenvalues in
\(\overline{\D}\).  Suppose \(H:\D\to M_n(\C)\) is analytic and
\[
 H(0)=I,\qquad \RePart H(w)\succeq0,
\qquad
 H(w)-F_\rho(w,T)\in\alg(T^*) .
\]
For \(0<t<1\), define
\[
 P_t=\sum_{m=0}^\infty t^{2m}T^{*m}T^m,\qquad
 Q_t=\sum_{m=0}^\infty t^mT^{*m}T^m,\qquad
 Y_t=Q_t-P_t.
\]
Then
\begin{equation}
\label{eq:rho-gramian}
 \frac2{1-t}Y_t-Y_tP_t-P_tY_t
 -(2-\rho)(P_t^2-P_t)\succeq0 .
\end{equation}
\end{proposition}

\begin{proof}
Diagonalize \(T=S\Lambda S^{-1}\), put \(G=S^*S\), and work first in
the eigenvector coordinates.  Polynomial interpolation identifies
\(\alg(T^*)\) with the diagonal algebra in these coordinates.  Thus
the transformed completion has the form
\[
 \widetilde H(w)
 =
 (1-k)G+kG(I-w\Lambda)^{-1}+\Theta(w)G,
 \qquad k=\frac2\rho,
\]
where \(\Theta(w)\) is diagonal analytic and \(\Theta(0)=0\).

The matrix Herglotz theorem implies that
\[
 \mathcal K_H(w,z)=\frac{H(w)+H(z)^*}{1-w\overline z}
\]
is a positive kernel.  Use its positivity at
\[
 w_i=t\overline{\lambda_i}\quad(1\leq i\leq n),
 \qquad w_0=0.
\]
In eigenvector coordinates let
\[
 P_{ij}=\frac{G_{ij}}{1-t^2\overline{\lambda_i}\lambda_j},
 \qquad
 Q_{ij}=\frac{G_{ij}}{1-t\overline{\lambda_i}\lambda_j},
 \qquad Y=Q-P.
\]
For an arbitrary \(u\in\C^n\), use the coordinate vectors
\(u_ie_i\) at \(w_i\) and the compensating vector
\[
 v=-G^{-1}Pu
\]
at the origin.  Exactly as in the unparameterized completion
argument, this choice cancels every value of the unknown diagonal
function \(\Theta\).  Explicitly, the entire correction contribution
to the kernel quadratic form is
\[
 2\RePart\sum_{i=1}^n
 \overline{u_i}\theta_i(w_i)
 \left((Gv)_i+(Pu)_i\right)=0.
\]

The remaining main block, cross block, and origin block are
respectively
\[
 2P+\frac{2k}{1-t}Y,\qquad
 (2-k)G+kQ,\qquad 2G.
\]
Here one uses the scalar partial-fraction identity
\[
 \frac1{(1-ta)(1-t^2a)}
 =
 \frac1{1-t}\left(\frac1{1-ta}
 -\frac{t}{1-t^2a}\right).
\]
Substitution of \(v=-G^{-1}Pu\) into the positive-kernel quadratic
form, followed by division by \(k\), gives
\[
 \frac2{1-t}Y-YG^{-1}P-PG^{-1}Y
 +(2-\rho)(P-PG^{-1}P)\succeq0.
\]
After congruence by \(G^{-1/2}\), put
\[
 \widetilde T=G^{1/2}\Lambda G^{-1/2}.
\]
The three balanced matrices are precisely the series in the
statement with \(T\) replaced by \(\widetilde T\).  The polar
decomposition \(S=UG^{1/2}\) gives
\(T=U\widetilde T U^*\), so conjugating back by \(U\) yields
\eqref{eq:rho-gramian} as stated.
\end{proof}

\begin{corollary}[The square-zero model of the Gramian inequality]
\label{cor:square-zero-completion}
Let \(T^2=0\), and suppose that \eqref{eq:rho-gramian} holds for
every \(0<t<1\).  Then
\[
 \|T\|\leq
 \begin{cases}
 \rho,&1\leq\rho\leq2,\\[1mm]
 \sqrt{\rho/(2-\rho)},&0<\rho\leq1.
 \end{cases}
\]
In particular, \(\|T\|\leq\max\{1,\rho\}\).
\end{corollary}

\begin{proof}
Put \(X=T^*T\).  Then
\[
 P_t=I+t^2X,\qquad Y_t=t(1-t)X.
\]
Equation \eqref{eq:rho-gramian} reduces to
\[
 \rho X-t(2-\rho t)X^2\succeq0.
\]
Consequently
\[
 \|T\|^2=\|X\|
 \leq\inf_{0<t<1}\frac{\rho}{t(2-\rho t)}.
\]
If \(\rho\geq1\), the denominator is maximized at \(t=1/\rho\),
giving \(\|T\|\leq\rho\).  If \(\rho\leq1\), its supremum is
\(2-\rho\), attained as \(t\uparrow1\), and the second bound follows.
\end{proof}

\begin{problem}[Parameterized completion theorem]
\label{prob:parameterized-completion}
Under the hypotheses of Proposition~\ref{prop:rho-gramian}, prove
\[
 \|T\|\leq\max\{1,\rho\}.
\]
\end{problem}

Corollary~\ref{cor:square-zero-completion} shows that the Gramian
inequality has at least the required strength on the square-zero
model, and is exactly sharp there when \(\rho\geq1\).
(The simple-spectrum hypothesis in
Proposition~\ref{prop:rho-gramian} is a device for deriving the
inequality, not an assumption in the corollary.)  For a general power sequence,
\eqref{eq:rho-gramian} alone has not yet yielded the asserted norm
bound; a proof must exploit either the full positive kernel or the
coupling of \eqref{eq:rho-gramian} over all \(t\).

\section{The exact geometric completion still required}

Let \(B\) have simple spectrum and let \(f\) be holomorphic on a
neighbourhood of \(\Omega_r(B)\), normalized so that
\(\|f\|_{\Omega_r(B)}\leq1\), with
\(\alg(f(B))=\alg(B)\), and put \(T=f(B)\).
The ordinary double-layer calculus produces a positive-real
completion with parameter \(2\).  To obtain the conjectured constant
one needs the following strengthened, algebra-aware statement.

\begin{problem}[Geometric \(q\)-completion]
\label{prob:geometric-completion}
With
\[
 \rho=\frac{2r}{1+\sqrt{1-r^2}},
\]
construct an analytic \(H:\D\to M_n(\C)\) such that
\[
 H(0)=I,\qquad \RePart H(w)\succeq0,\qquad
 H(w)-F_\rho(w,T)\in\alg(T^*).
\]
\end{problem}

A sufficient, stronger formulation is to construct a positive operator-valued
boundary measure \(M\) on \(\partial\Omega_r(B)\) having mass
\(\rho I\) and satisfying, for every holomorphic \(h\),
\[
 \int h\,dM-h(B)\in\alg(B^*).
\]
Indeed, in that case
\[
 H(w)=\frac1\rho
 \int_{\partial\Omega_r(B)}
 \frac{1+w f(z)}{1-w f(z)}\,dM(z)
\]
has positive real part and \(H(0)=I\).  Expanding the Cayley kernel
in powers of \(w\), and using
\(\alg(B^*)=\alg(f(B)^*)\), gives
\[
 H(w)-F_\rho(w,f(B))\in\alg(f(B)^*).
\]

\begin{proposition}[Exact obstruction to scalar stripping]
\label{prop:scalar-stripping}
Let \(0<r<1\) and \(A=\diag(-1,1)\).  If one subtracts the pointwise
smallest eigenvalue from the standard double-layer measure on
\(\partial\Omega_r(A)\), the remaining mass is
\[
 \rho_{\mathrm{strip}}(r)I,\qquad
 \rho_{\mathrm{strip}}(r)
 =2-\frac4\pi\arccos r.
\]
This is strictly larger than the desired parameter
\[
 \rho_r=\frac{2r}{1+\sqrt{1-r^2}}.
\]
Thus pointwise scalar stripping cannot produce the measure required
in Problem~\ref{prob:geometric-completion}, even for this normal
\(2\times2\) matrix.
\end{proposition}

\begin{proof}
Put \(s=\sqrt{1-r^2}\).  The scaled \(q\)-range is the ellipse
parametrized by
\[
 \zeta(\phi)=\frac{\cos\phi+is\sin\phi}{r},
 \qquad 0\leq\phi\leq2\pi.
\]
Direct substitution in the scalar double-layer formula shows that
the two diagonal densities, with respect to \(d\phi\), are
\[
 m_\lambda(\phi)
 =\frac1\pi\RePart\frac{-i\zeta'(\phi)}
 {\zeta(\phi)-\lambda},
 \qquad
 m_+(\phi)=\frac{s}{\pi(1-r\cos\phi)},
 \qquad
 m_-(\phi)=\frac{s}{\pi(1+r\cos\phi)}.
\]
Each has integral \(2\).  The scalar mass that can be removed while
preserving pointwise positivity is therefore
\begin{align*}
 a(r)
 &=\int_0^{2\pi}\min\{m_+(\phi),m_-(\phi)\}\,d\phi\\
 &=\frac{s}{\pi}\int_0^{2\pi}
   \frac{d\phi}{1+r|\cos\phi|}
 =\frac4\pi\arccos r.
\end{align*}
The remaining mass is \(2-a(r)\), proving the first formula.

Write \(r=\sin\theta\), \(0<\theta<\pi/2\).  Then
\[
 \rho_{\mathrm{strip}}(r)=\frac{4\theta}{\pi},
 \qquad
 \rho_r=2\tan\frac{\theta}{2}.
\]
The strict convexity of \(\tan x\) on \([0,\pi/4]\) places its graph
strictly below the chord joining \((0,0)\) and \((\pi/4,1)\).
Consequently
\[
 \tan\frac{\theta}{2}<\frac{2\theta}{\pi},
\]
which proves \(\rho_{\mathrm{strip}}(r)>\rho_r\).
\end{proof}

The ordinary double-layer measure has mass \(2I\).  Thus the issue is
not merely to preserve positivity, but to lower the mass to \(\rho I\)
while retaining all of the stated holomorphic moments modulo
\(\alg(B^*)\).  Proposition~\ref{prop:scalar-stripping} shows that
scalar subtraction does not suffice; the adjoint-algebra freedom is
the additional structure available for a successful construction.

\begin{proposition}[Conditional closure of the conjecture]
\label{prop:conditional-closure}
If Problems~\ref{prob:parameterized-completion} and
\ref{prob:geometric-completion} hold, then
Problem~\ref{prob:q-conjecture} holds for every finite-dimensional
matrix.
\end{proposition}

\begin{proof}
For simple spectrum and an algebra-generating auxiliary \(f\), the geometric
completion supplies the hypotheses of the parameterized completion
theorem, hence
\[
 \|f(B)\|\leq\max\{1,\rho\}=C_r.
\]
Let \(p\) be an arbitrary polynomial.  If
\(\lambda_1,\ldots,\lambda_n\) are the distinct eigenvalues of \(B\),
then
\[
 p_\varepsilon(z)=p(z)+\varepsilon z
\]
takes distinct values on the \(\lambda_i\)'s for all but finitely many
\(\varepsilon\).  Hence
\(\alg(p_\varepsilon(B))=\alg(B)\).  Apply the preceding estimate to
\[
 f_\varepsilon=
 \frac{p_\varepsilon}
 {\|p_\varepsilon\|_{\Omega_r(B)}}
\]
and let \(\varepsilon\to0\).  This proves the conjectured estimate for
every polynomial when \(B\) has simple spectrum.

Finally choose matrices \(B_j\to A\) with simple spectra.  The same
compact parameter space of admissible pairs in the definitions gives
\[
 d_H\bigl(W_r(B_j),W_r(A)\bigr)\leq\|B_j-A\|,
\qquad
 d_H\bigl(\Omega_r(B_j),\Omega_r(A)\bigr)
\leq\frac{\|B_j-A\|}{r}.
\]
Therefore \(p(B_j)\to p(A)\), while
\[
 \sup_{\Omega_r(B_j)}|p|
 \longrightarrow
 \sup_{\Omega_r(A)}|p|.
\]
Passing to the limit proves the result for \(A\).
\end{proof}

\section{Conclusion}

The sharp constant and its phase transition are now proved for all
quadratic matrices and their reducing direct sums.  The
two-dimensional proof is not merely a check of a special example:
the exact extremizers \(I\) and \(J\) already force both branches of
the conjectured constant.  In the general case the argument has been
reduced to Problems~\ref{prob:parameterized-completion} and
\ref{prob:geometric-completion}.  The first is an operator-theoretic
completion question; the second is a geometric moment problem on the
boundary of the scaled \(q\)-numerical range.  Closing only one of
them is insufficient, and the scalar-floor version of the second
loses essential algebraic information.

\begin{thebibliography}{99}

\bibitem{ChienNakazatoPsarrakos2005}
M.-T. Chien, H. Nakazato, and P. Psarrakos,
\emph{On the \(q\)-numerical range of matrices and matrix
polynomials},
Linear Multilinear Algebra \textbf{53} (2005), no.~5, 357--374.
\url{https://doi.org/10.1080/03081080500167596}.

\bibitem{DLMF}
NIST Digital Library of Mathematical Functions,
Chapters 20 and 22.
\url{https://dlmf.nist.gov/}.

\bibitem{KanasSugawa2006}
S. Kanas and T. Sugawa,
\emph{On conformal representations of the interior of an ellipse},
Ann. Acad. Sci. Fenn. Math. \textbf{31} (2006), no.~2, 329--348.

\bibitem{LiPoonSze2011}
C.-K. Li, Y.-T. Poon, and N.-S. Sze,
\emph{Elliptical range theorems for generalized numerical ranges of
quadratic operators},
Rocky Mountain J. Math. \textbf{41} (2011), no.~3, 813--832.
\url{https://doi.org/10.1216/RMJ-2011-41-3-813}.

\bibitem{LiMehtaRodman1994}
C.-K. Li, P.~P. Mehta, and L. Rodman,
\emph{A generalized numerical range: the range of a constrained
sesquilinear form},
Linear Multilinear Algebra \textbf{37} (1994), no.~1--3, 25--49.
\url{https://doi.org/10.1080/03081089408818311}.

\bibitem{OLoughlinRani2026}
R. O'Loughlin and J. Rani,
\emph{\(q\)-Numerical ranges and spectral sets},
arXiv:2603.15536, 2026.
\url{https://arxiv.org/abs/2603.15536}.

\bibitem{PagaczPietrzyckiWojtylak2020}
P. Pagacz, P. Pietrzycki, and M. Wojtylak,
\emph{Between the von Neumann inequality and the Crouzeix
conjecture},
Linear Algebra Appl. \textbf{605} (2020), 130--157.
\url{https://doi.org/10.1016/j.laa.2020.07.002}.

\end{thebibliography}

\end{document}
